\SourceLeaf{131}
Resultate zu gewinnen, sondern nur allgemeine Ueberschläge und Anhaltspunkte, so ist es doch wichtig, diese mit den Resultaten einer strengern Methode in Einklang zu finden, und jedenfalls wird dadurch alles in ein helleres Licht gesetzt. Zudem sind diese Auseinandersetzungen enge verknüpft mit der prüfenden Revision eines Cardinalpunkts des jetzt bestehenden Regulativs, welche Prüfung ich für unumgänglich nothwendig halte, und in Beziehung auf welche keine Dunkelheiten zurückbleiben dürfen. Ich habe daher diese Entwicklungen in dem ersten Theile dieses Aufsatzes so ausführlich und, wie ich hoffe, so klar dargestellt, dass man denselben leicht wird folgen können.

\Needspace{5\baselineskip}\pdfbookmark[2]{Erste Abtheilung}{witwen-II-1}
\begin{center}\textit{Erste Abtheilung.}\end{center}
\noindent\hspace*{\parindent}Dass der Zustand der Witwenkasse bei dem Senate zur Sprache gebracht ist, und Verhandlungen darüber Statt gefunden haben, in deren Folge eine gründliche Untersuchung jenes Zustandes von dem Universitäts-Curatorium verfügt ist, war zunächst durch die im Herbst des vorigen Jahrs hervorgetretenen Besorgnisse veranlasst, welche besonders durch das rasche und alle frühern Erfahrungen weit überschreitende Steigen der Witwenzahl (seit dem Tode des Geheimen Justizraths M. auf zwei und zwanzig) erregt, und durch eine augenblickliche Insufficienz des baaren Kassenvorraths zur vollständigen Zahlung der Pensionen am gewohnten Tage noch vergrössert waren.

Dass diese und andere Umstände eine gewisse Beunruhigung hervorbrachten, ist um so weniger zu verwundern, da man sich gewöhnt hatte, den Zustand wie einen höchst blühenden zu betrachten. Bis Ostern 1829 war der Betrag der jährlichen Pension 210 Thl. gewesen, und durch viermalige Erhöhung von je 10 Thl. während des kurzen Zeitraums von \(6\tfrac12\) Jahren war sie Michaelis 1835 auf 250 Thl. gestiegen. Man glaubte daher damals den Zeitpunkt, wo die Pensionen auf 300 Thl. angewachsen sein würden, so nahe, dass man sich schon mit Planen beschäftigte, wie nachher der Ueberfluss am besten zu verwenden sein würde\footnote{Zwei wohlmeinende, seitdem bereits verstorbene Mitglieder der Kirchen-Deputation brachten in Anregung, der eine die Abschaffung der jährlichen Beiträge, der andere die Erweiterung der Waisenpensionen, bis zur Stiftung lebenslänglicher Pensionen für die unverheiratheten Professorentöchter.}. Allein gerade von jener Zeit an begannen die Verhältnisse sich ungünstiger zu gestalten; zu den bis Ende 1835 vorhandenen zwölf Pensionirten kamen binnen 9 Jahren zwölf neue Witwen hinzu, während nur zwei Pensionen erloschen.

Die vorhin erwähnte augenblickliche Unzulänglichkeit des baaren Kassenbestandes ist übrigens ein Umstand von geringer Bedeutung, selbst wenn dadurch eine kurzfristige verzinsliche Anleihe nöthig geworden wäre, was jedoch, der Jahresrechnung 1844—1845 zufolge, dasmal nicht der Fall gewesen zu sein scheint. Dergleichen Eventualitäten können bei der solidesten Kasse, wie bei dem solidesten in vielfachem Geldverkehr begriffenen Particulier vorkommen, und desto öfter, je mehr dahin gestrebt wird, grössere Geldsummen nicht lange ungenutzt liegen zu lassen.

Auch die 1844 so sehr vergrösserte Anzahl der Witwen war, an sich betrachtet, noch kein Beweis einer nahen Gefahr. Man durfte mit Wahrscheinlichkeit erwarten, dass diese Zahl bald wieder eine Verminderung erleiden würde, wie denn auch wirklich von März bis Juni d. J. drei Witwen mit Tode abgegangen sind. Auch ist nicht unwahrscheinlich, dass in nicht langer Zeit noch einige weitere Abnahme eintreten könne: indessen gewährt eine Rechnung von heute auf morgen nur eine sehr ungenügende Beruhigung, und ein beträchtliches dauerndes Sinken der Witwenzahl hat man allerdings keinen Grund zu erwarten.

Nicht die zeitweilige Grösse der Witwenzahl ist es, was dem Institute Gefahr drohet, sondern etwas ganz anderes, nemlich
